\documentclass[10pt,a4paper]{amsart}
\usepackage[margin=19mm]{geometry}
\usepackage{amsmath,amssymb,mathtools}
\usepackage[scr=rsfs,cal=euler]{mathalfa}
\usepackage{xfrac}
\usepackage[unicode,hidelinks,bookmarks=false]{hyperref}
\newcommand{\ie}{i.e.\ }
\newcommand{\cf}{cf.\ }
\newcommand{\resp}{respectively\ }
\newcommand{\Subsection}{\subsection}
\providecommand{\nobreakdash}{\nobreak\hbox{-}}
\setlength{\emergencystretch}{2em}
\multlinegap0pt
\usepackage{tikz}
\usepackage{algorithm}\usepackage{algpseudocode}
\usepackage{enumerate}
\usepackage{mathtools}
\usepackage{amssymb}
\newtheorem{lemma}{Lemma}
\title{On the Kolmogorov-Feller weak law of large numbers for the Fr\'{e}chet mean on non-compact symmetric spaces}
\author{Jongmin Lee, Sungkyu Jung}
\date{Source: arXiv:2509.02074v2 (CC BY 4.0)}
\begin{document}
\maketitle

\noindent\emph{Source and licence.} On the Kolmogorov-Feller weak law of large numbers for the Fr\'{e}chet mean on non-compact symmetric spaces. Version arXiv:2509.02074v2 (CC BY 4.0), distributed by arXiv under the Creative Commons Attribution 4.0 International licence, http://creativecommons.org/licenses/by/4.0/, as recorded in the arXiv licence metadata for this exact version and shown on the abstract page https://arxiv.org/abs/2509.02074v2. Full licence notice: \texttt{licenses/arxiv-2509.02074-CC-BY-4.0.txt}.\par\smallskip

\noindent\textit{Editorial scope.} Counted unit: the article's lemma lem:somevsall\_tail\_conditions (source0.tex lines 703--715). The article's complete original proof of it is delivered with it (source0.tex lines 718--785, one \texttt{proof} environment of 68 lines). No mathematics is written, repaired or re-derived: the statement and the proof are the article's own lines, in the article's own notation, and no retained line is substituted or re-flowed.  Retained uncounted as the source-local context the proof needs: lines 670--675.  One declared adaptation: exactly 1 ancillary strings inside the retained spans are omitted (image-inclusion commands and, where the source repeats a label, the repeated label definitions) and no other line differs from the source; the figure environments, with their placement, captions and labels, are kept, and the package records the omitted strings in fidelity receipts bound to the affected bodies.  No retained line contains a citation, so the wrapper carries no bibliography and no bibliography entry.  The retained lines contain the article's own diagram code, which the wrapper typesets with the standard tikz packages; no image file is delivered and the package declares no asset dependency.  Editorial formatting only, in addition: an AMS \texttt{amsart} preamble that restates the article's own declarations and macros used by the retained lines, namely the environments \texttt{lemma} on separate counters, together with the standard packages those lines need.  The retained environments print in this document as Lemma 1 in order of appearance, whereas the article prints the same items with the numbering of its own full document.  The complete native source of the article as distributed in the arXiv e-print for this exact version is bundled unchanged as \texttt{sources/184299/source0.tex}.
\begin{figure}[htbp!]
\centering

\caption{An illustration of the transvection $T$, mapping from $x$ to $s_{\mu_2}\circ s_{\mu_1}(x)$.}
\label{fig:transvection}
\end{figure}

\begin{lemma}\label{lem:somevsall_tail_conditions} 
Let $(M,d)$ be a metric space, and let $(X_i)_{i=1}^\infty$ be a sequence of independent $M$-valued random variables. For an $m \in M$, let $c_1(m)$, $c_2(m)$ and $c_3(m)$ denote the following conditions, respectively: 
\begin{align*}
   c_1(m): & \lim_{n\to\infty} n\mathbb{P}(d(X_1,m) > n) = 0;\\
   c_2(m): & \lim_{n\to\infty} n^{-2}\sum_{i=1}^n\mathbb{E}[d^2\{\overline{X_{n,i}}(m),m\}] = 0;\\
   c_3(m): & \lim_{n\to\infty} \sum_{i=1}^n\mathbb{P}(d(X_i,m) > n) = 0.
\end{align*}
\begin{itemize}
    \item[(i)] For any $m,m' \in M$, $c_1(m)$ holds if and only if $c_1(m')$ holds.
    \item[(ii)] For any $m,m' \in M$, $c_2(m)$ holds if and only if $c_2(m')$ holds.
    \item[(iii)] Additionally assume that $M$ is a symmetric space and,  for a given $\mu \in M$, $\mathbb{P}_{X_i}$ is geodesically symmetric about $\mu$. Then, for any $m \in M$, $c_3(\mu)$ holds if $c_3(m)$ holds.
\end{itemize}
\end{lemma}

\begin{proof}[Proof of Lemma \ref{lem:somevsall_tail_conditions}]
(i) Let $m, m'$ be fixed, and without losing generality, assume that $R:= d(m,m') > 0$. We shall show that the condition $c_1(m)$ implies $c_1(m')$. The converse is obtained by the symmetry of the argument. 

Let $x \in M$ be any point. 
By the triangle inequality, 
\begin{equation}
   d(x,m') \le d(x, m) + d(m,m') = d(x, m) + R. \label{eq:triangle}
\end{equation} 
Therefore, for any $n > R$, $d(x,m') > n$ implies $d(x,m)> n - R$, which in turn implies that 
$$\{x \in M: d(x,m') > n\} \subset \{x \in M: d(x,m)> n - R\}.$$
Then, 
\begin{align*}
    n\mathbb{P}(d(X_1,m')>n) & \le n\mathbb{P}(d(X_1,m)>n-R) \\
      & = \frac{n}{n-R} (n-R) \mathbb{P}(d(X_1,m)>n-R).
\end{align*}
But, since we assumed $c_1(m)$, we have $(n-R) \mathbb{P}(d(X_1,m)>n-R) \to 0$ as $n\to\infty$, and $\lim_{n\to\infty}{n}/({n-R})= 1$ as $R$ is fixed. Thus, $n\mathbb{P}(d(X_1,m')>n) \to 0$ as $n\to\infty$, which is precisely the condition $c_1(m')$. 

(ii) Let $m, m'$ be fixed, and let $R:= \lceil d(m,m') \rceil > 0$ be the smallest non-zero integer larger than or equal to $d(m,m')$. As in Part (i), it suffices to show that $c_2(m)$ implies $c_2(m')$. 
Observe first that, for any $X_i,m \in M$, 
$$d^2(\overline{X_{n,i}}(m),m) = d^2(X_i,m) 1_{d(X_i,m) \le n},$$
which can be deduced from the definition of $\overline{X_{n,i}}(m)$. 

By the triangle inequality %(\ref{eq:triangle})
we have 
$$\{x \in M: d(x,m') \le n\} \subset \{x \in M: d(x,m)\le n + R\},$$ 
and, for any $x \in M$ 
$$d^2(x,m') \le \{d(x,m) + R\}^2 \le 2d^2(x,m) + 2R^2.$$
It follows that
\begin{align*}
    d^2(X_i,m') 1_{d(X_i,m') \le n} 
    & \le [2d^2(X_i,m) + 2R^2]1_{d(X_i,m') \le n} \\ 
    & \le [2d^2(X_i,m) + 2R^2]1_{d(X_i,m)\le n + R}.
\end{align*}

Write $F_m(n) : = n^{-2}\sum_{i=1}^n\mathbb{E}[d^2\{\overline{X_{n,i}}(m),m\}]$ for any $m\in M$ and $n\ge 1$. Then, 
\begin{align*}
    F_{m'}(n) & = n^{-2}\sum_{i=1}^n\mathbb{E} d^2(X_i,m') 1_{d(X_i,m') \le n}  \\
    & \le \frac{2}{n^2}\sum_{i=1}^n\mathbb{E} d^2(X_i,m)1_{d(X_i,m)\le n + R} + \frac{2R^2}{n^2}\sum_{i=1}^n \mathbb{P}(d(X_i,m)\le n + R) \\ 
    & \le \frac{2(n+R)^2}{n^2}F_{m}(n+R) + \frac{2R^2}{n^2}\sum_{i=1}^n 1.
\end{align*}
Since $F_{m}(n+R) \to 0$ as $n\to \infty$ (which is precisely the condition $c_2(m)$), $F_{m'}(n) \underset{n\to \infty}{\to} 0$ as desired. 





(iii) Recall that $s_\mu(\cdot)$ is the geodesic symmetry with respect to $\mu \in M$. For any $x \in M$, we have 
$d(x, s_\mu(x)) = 2d(x,\mu) = 2d(s_\mu(x),\mu)$ (see Figure~\ref{fig:transvection} for an illustration). By the triangle inequality, for any $m\in M$, we obtain 
\begin{equation}
    \label{eq:iii} 
    2d(x,\mu) = d(x, s_\mu(x)) \le d(x,m ) + d(m, s_\mu(x)).
\end{equation}

Let us now fix $m$, and suppose that $c_3(m)$ holds. 
From (\ref{eq:iii}), we have 
\begin{align*}
\{x\in M: d(x,\mu) > n\} &\subset \{x\in M: d(x,m) + d(s_\mu(x), m) > 2n\}\\
&\subset \{x\in M: d(x,m) > n\}\cup \{x\in M: d(s_\mu(x),m) > n\}.
\end{align*}
Since $\mathbb{P}_{X_i}$ is geodesically symmetric about $\mu$, $X_i$ and $s_\mu (X_i)$ have the same distribution. Therefore, we have 
$$
\mathbb{P}(d(X_i,\mu) > n ) \le 2\mathbb{P}(d(X_i,m) > n).$$
Summing over $1\le i\le n$, we have 
$$
\sum_{i=1}^n\mathbb{P}(d(X_i,\mu) > n ) \le 2\sum_{i=1}^n\mathbb{P}(d(X_i,m) > n) \to 0,
$$
as $n\to \infty$. That is, $c_3(\mu)$ holds.
\end{proof}

\end{document}
